\documentclass[../../../include/open-logic-section]{subfiles}

\begin{document}

\olfileid{sth}{story}{urelements}
\olsection{મૂળઘટકો સ્વીકારવા કે નહીં?}

આગળનાં થોડાં પ્રકરણોમાં આપણે ગણની સંચયી-પુનરાવર્તી કલ્પનામાંથી
સ્વયંસિદ્ધિઓ મેળવવાનો પ્રયાસ કરીશું. પરંતુ આગળ વધતાં પહેલાં
\emph{મૂળઘટકો} વિશે થોડું વધુ કહેવું જરૂરી છે.

\olref[approach]{sec}ના ચિત્રમાં જૂના \emph{ગણો}માંથી જ નવા ગણો
રચી શકતા હતા. પરંતુ કદાચ આપણે અમુક \emph{ગણ ન હોય તેવી વસ્તુઓ}---
ગાયો, ડુક્કરો, રેતીના કણો કે બીજું કંઈ---ગણોના !!{element}s
હોઈ શકે એવું સ્વીકારવા ઇચ્છીએ. તે કિસ્સામાં અમુક પાયાના ઘટકો,
\emph{મૂળઘટકો},થી શરૂ કરીએ અને પછી કહીએ કે દરેક તબક્કા $S$ પર
મૂળઘટકો અથવા $S$ પહેલાંના તબક્કાઓમાં રચાયેલા ગણોથી બનેલા
``બધા શક્ય'' ગણો (કોઈ પણ મિશ્રણમાં) રચીએ છીએ.
મળતું ચિત્ર નીચેના જેવું વધુ લાગે:
\begin{center}
\begin{tikzpicture}[scale=0.6]
	\tikzset{cut_here/.style={densely dotted}}
	\draw (-4,6) -- (-1, 0) -- (1,0) -- (4,6); 
	\draw[cut_here] (-5,8)--(-4,6);
	\draw[cut_here] (5,8)--(4,6);
	\draw(-1.5, 1)--(1.5, 1);
	\draw(-2, 2)--(2, 2);
	\draw(-2.5, 3)--(2.5,3);
	\draw(-3, 4)--(3, 4);
	\draw(-3.5, 5)--(3.5, 5);
	\draw(-4, 6)--(4, 6);
	\node[label] at (2, 0) {\small 0};
	\node[label] at (2.5, 1) {\small 1};
	\node[label] at (3, 2) {\small 2};
	\node[label] at (3.5, 3) {\small 3};
	\node[label] at (4, 4) {\small 4};
	\node[label] at (4.5, 5) {\small 5};
	\node[label] at (5, 6) {\small 6};
	\end{tikzpicture}
\end{center}
તેથી હવે નિર્ણય લેવાનો છે: \emph{શું મૂળઘટકો સ્વીકારવા જોઈએ?}

દાર્શનિક રીતે આપણા સિદ્ધાંતમાં મૂળઘટકો સામેલ કરવા અર્થપૂર્ણ છે.
મુખ્ય કારણ ગણસિદ્ધાંતને \emph{લાગુ કરી શકાય એવો} બનાવવાનું છે.
મુદ્દો સમજવા માટે \olref[sfr][siz][]{chap}માંથી યાદ કરો કે બે ગણો
$A$ અને~$B$નું કદ સમાન છે, એટલે કે $\cardeq{A}{B}$, એમ ત્યારે
અને માત્ર ત્યારે કહીએ છીએ જ્યારે તેમની વચ્ચે દ્વિએક વિધેય હોય.
હવે જો ખેતરની ગાયો અને વાડાના ડુક્કરો બંને ગણો બનાવતા હોય, તો
``ગાયો જેટલાં જ ડુક્કરો છે'' એવા દાવાનું ગણસિદ્ધાંત મુજબ વિશ્લેષણ
કરી શકીએ. પરંતુ જો મૂળઘટકોને પ્રતિબંધિત કરીએ અને તેથી ગાયો તથા
ડુક્કરો ગણો બનાવતાં \emph{નથી}, તો આવું વિશ્લેષણ ઉપલબ્ધ નહીં હોય.
ખરેખર, ગણો પોતે સિવાયની કોઈ પણ વસ્તુને ગણસિદ્ધાંત લાગુ કરવાની
સીધી રીત આપણી પાસે નહીં હોય. (મૂળઘટકો સામેલ કરવાનાં વધુ કારણો
માટે \citealt[pp.~vi, 24, 50--1]{Potter2004} જુઓ.)

જોકે ગણિતમાં મૂળઘટકો સ્વીકારવાનું બહુ ઓછું જોવા મળે છે. તેનું
એક કારણ એ છે કે મૂળઘટકો વિના ગણસિદ્ધાંત ઘડવો \emph{ખૂબ જ થોડુંક}
સરળ છે. પરંતુ ક્યારેક તેમને ગણસિદ્ધાંતમાંથી બાકાત રાખવા માટે
વધુ રસપ્રદ કારણો મળે છે:
\begin{quote}
ગણસિદ્ધાંત ગણિતનો પાયો છે એવી માન્યતા મુજબ માત્ર ગણો વિશે વાત
કરીને આખું ગણિત વ્યક્ત કરી શકવું જોઈએ. તેથી આપણા ચલોએ ગાયો
અને ડુક્કરો જેવી વસ્તુઓ પર ચલવું ન જોઈએ.
%But if $C$ is a cow, $\{C\}$ is a set, but not a legitimate mathematical object.
\citep[p.~8]{Kunen1980}
\end{quote}
આમ લાગુ કરી શકવાની ક્ષમતા પર ભાર મૂકવાથી મૂળઘટકોને
\emph{સામેલ કરવા} સૂચવાય; ગણિતને શુદ્ધ ગણસિદ્ધાંત સુધી ઘટાડવાના
પાયાલક્ષી હેતુ પર ભાર મૂકવાથી તેમને \emph{બાકાત રાખવા} સૂચવાય.
થોડી આળસ પણ મૂળઘટકો બાકાત રાખવાની દિશા તરફ દોરી જાય છે.

આપણે સૌથી આળસુ માર્ગ લઈશું. જોકે તેની પાછળ થોડું શિક્ષણલક્ષી
કારણ પણ છે. આપણો હેતુ ગણસિદ્ધાંતના દર્શનશાસ્ત્રનો અભ્યાસ શરૂ
કરવા માટે જરૂરી ગણસિદ્ધાંતની પાયાની બાબતોનો પરિચય આપવાનો છે.
અને તે દાર્શનિક સાહિત્યનો મોટા ભાગનો હિસ્સો મૂળઘટકો
\emph{વિના} ઘડાયેલા ગણસિદ્ધાંતોની ચર્ચા કરે છે. તેથી આ પુસ્તક
તે સાહિત્યને નજીકથી અનુસરે તો કદાચ વધુ ઉપયોગી બને.

\end{document}
